\documentclass[12pt]{article}
\usepackage[margin=1in]{geometry}
\usepackage{amsmath}
\usepackage{natbib}
\usepackage{hyperref}

\title{State Defaults and Constitutional Reform in the 1840s:\\
Pennsylvania and New York's Response to the Debt Crisis}
\author{Claud and Tom}
\date{\today}

\begin{document}

\maketitle

\begin{abstract}
In the early 1840s, multiple U.S. states defaulted on bonds issued to finance infrastructure projects, precipitating a wave of constitutional reforms that fundamentally reshaped state fiscal institutions. This paper examines the defaults of Pennsylvania and New York's near-default, and the subsequent constitutional restrictions on state borrowing adopted by both states. These reforms, which required voter approval of taxes before debt could be incurred and prohibited state involvement in private corporations, effectively eliminated ``taxless finance'' and established fiscal constraints that remain influential in American public finance today.
\end{abstract}

\section{Introduction}

In the early 1840s, the United States experienced a sovereign debt crisis at the state level. Nine states and one territory defaulted on bond obligations following the Panic of 1837, marking one of the most significant fiscal crises in American history. This crisis led to fundamental constitutional reforms in multiple states, with New York and Pennsylvania adopting strict limitations on state borrowing and prohibitions on state credit being pledged to private enterprises.

The defaults were concentrated among states that had borrowed heavily to finance ``internal improvements''---primarily canals and railroads---using what economic historians term ``taxless finance.'' Rather than raising property taxes to fund infrastructure, states borrowed against anticipated toll revenues and hoped that increased land values would eventually generate sufficient tax receipts to service the debt \citep{wallis2005constitutions}.

\section{The Default Crisis}

\subsection{Pennsylvania's Default}

Pennsylvania defaulted in 1842 after accumulating over \$36 million in debt (approximately \$1.1 billion in 2025 dollars), primarily to finance its ``Main Line'' canal system connecting Philadelphia to Pittsburgh and other internal improvements \citep{wallis2004sovereign}. The state's debt burden was the third-largest in absolute terms, behind only New York and Maryland.

Pennsylvania pursued an aggressive strategy of infrastructure investment without concurrent tax increases. The state relied on several expedients to postpone fiscal reckoning:
\begin{itemize}
\item In 1836, the state chartered the Bank of the United States of Pennsylvania (BUSP), which paid the state over \$12 million in bonus payments, loans, and investments
\item In 1837, the distribution of the federal surplus provided Pennsylvania with \$2.87 million
\item The state continued borrowing almost \$4 million more in 1840 even as financial conditions deteriorated
\end{itemize}

These temporary windfalls masked the fundamental problem: Pennsylvania had borrowed extensively while maintaining an aversion to direct taxation. As one historian observed, ``Speculation and hatred of all forms of direct taxation were the causes of the downfall in Pennsylvania's credit'' \citep{wallis2004sovereign}.

\subsection{New York's Near-Default}

New York narrowly avoided default despite having the largest state debt of \$22 million. The state's position was saved primarily by the extraordinary success of the Erie Canal, completed in 1825, which generated substantial toll revenues. However, the near-default experience and the fiscal pressure of the early 1840s motivated significant constitutional reform.

As \citet{wallis2005constitutions} documents, nine of the ten states with the largest per capita debts defaulted, with New York, Ohio, and Alabama narrowly avoiding this fate. The pattern was clear: states that had pursued large-scale infrastructure investment through debt financing, without corresponding tax increases, were most vulnerable to default when the economic depression following the Panic of 1837 reduced both toll revenues and land values.

\subsection{Other Defaulting States}

Beyond Pennsylvania, eight other states and the Florida Territory defaulted between 1841 and 1843:
\begin{itemize}
\item \textbf{Maryland}: Defaulted with over \$12 million in debt, primarily for canal construction
\item \textbf{Midwest states} (Indiana, Illinois, Michigan): Combined debt of approximately \$37 million for railroads and canals
\item \textbf{Southern states} (Arkansas, Louisiana, Mississippi): Borrowed primarily to capitalize state banks
\end{itemize}

Four states ultimately repudiated all or part of their debts, while three underwent substantial renegotiations \citep{grinath1997debt}.

\section{Constitutional Responses}

State legislatures throughout the country asked ``how did we get into this mess?'' and ``how can we prevent this from happening again?'' Although conditions in every state were unique, the constitutional reforms of the 1840s shared a common theme: states got into trouble because they pursued taxless finance, and the solution was to eliminate this option by requiring that taxes be raised when spending is contemplated \citep{wallis2005constitutions}.

\subsection{New York Constitution of 1846}

New York held a constitutional convention in 1846 that resulted in strict debt limitations. The key provisions, contained in Article VII of the new constitution, established what scholars call ``procedural debt restrictions'':

\begin{enumerate}
\item A debt limit of \$1 million for ordinary purposes (Section 1)
\item For debt exceeding this limit (Section 3), the legislature must:
\begin{itemize}
\item Contract the debt for a specific, stated purpose
\item Enact a tax sufficient to pay the interest and retire the principal within 18 years
\item Submit the proposed tax to voters for approval in a statewide referendum
\end{itemize}
\item Exceptions only for extraordinary emergencies: invasion, insurrection, or war
\end{enumerate}

This framework was revolutionary in American public finance. As described in the Harvard Business School case study on the convention, the proposal by Michael Hoffman would place ``tight restrictions on state debt, which had increased sharply over the previous eight years'' \citep{moss2016debt}. By requiring voter approval of taxes \emph{before} projects were undertaken, rather than after debt had accumulated, the constitution fundamentally changed the political calculus of infrastructure investment.

The procedural restrictions did not limit the \emph{amount} of debt a state could issue. Rather, by addressing the timing issue, they created what modern public finance literature would call a hard budget constraint. They made it impossible for project promoters to promise valuable infrastructure while deferring the tax burden to the future \citep{grinath1997debt}.

\subsection{Pennsylvania Constitutional Amendments of 1857}

Pennsylvania's constitutional response was more delayed than New York's. The state held a constitutional convention in 1837--38 that produced a new constitution adopted in 1838. However, this constitution did not include comprehensive debt restrictions. The major fiscal reforms came as amendments added in 1857, fifteen years after the default.

Article XI of the amended Pennsylvania Constitution established strict limitations:

\begin{itemize}
\item \textbf{Section 1}: Limited state debt to \$750,000 for ``casual deficits or failures in revenues''
\item \textbf{Section 3}: ``Except the debts above specified...no debt whatever shall be created by, or on behalf of the state''
\item \textbf{Section 4}: Created a mandatory sinking fund of at least \$250,000 annually to retire existing debt
\item \textbf{Section 5}: ``The credit of the commonwealth shall not in any manner, or event, be pledged, or loaned, to any individual, company, corporation, or association''
\item \textbf{Section 6}: Prohibited the state from assuming debts of local governments or corporations
\item \textbf{Section 7}: Prohibited counties and cities from becoming stockholders in corporations or lending their credit
\end{itemize}

These provisions effectively prohibited Pennsylvania from the practices that had led to default: borrowing without dedicated revenue sources, investing in or lending to private corporations, and assuming the debts of other entities.

\section{Common Elements of the Reforms}

The constitutional reforms in New York, Pennsylvania, and other states shared several common features that together eliminated the possibility of taxless finance:

\subsection{Strict Debt Limits}

Both states imposed absolute limits on state borrowing for ordinary purposes. New York's \$1 million limit and Pennsylvania's \$750,000 limit were intentionally restrictive, essentially requiring that any significant borrowing meet extraordinary criteria or undergo special approval processes.

\subsection{Prohibition on Lending State Credit}

A critical component was prohibiting the state from pledging its credit to, or becoming a stockholder in, private corporations. This addressed a major channel through which states had accumulated debt: making loans to railroad and canal companies, or purchasing stock in such enterprises with borrowed funds.

\subsection{Dedicated Revenue Requirements}

The reforms required that any authorized debt be accompanied by specific, dedicated tax revenues sufficient to service and retire it. This forced legislators and voters to confront the true cost of projects at the time of authorization, rather than deferring the tax burden.

\subsection{Voter Approval Requirements}

New York's requirement for voter approval of debt-servicing taxes was particularly innovative. By requiring a statewide referendum on the taxes needed to service debt, the constitution ensured that the decision to borrow would be made with full awareness of the fiscal consequences.

\section{The ``Taxless Finance'' Problem}

\citet{wallis2005constitutions} argues that the fundamental problem was not the infrastructure investments themselves---many projects, particularly New York's Erie Canal, were economically sound and generated substantial social returns. The problem was the political economy of deferred taxation.

When project promoters could promise infrastructure while claiming it would ``pay for itself'' through tolls or increased land values, voters and legislators systematically underestimated risks. Without the requirement to raise current taxes, several problems emerged:

\begin{enumerate}
\item \textbf{Optimism bias}: Toll revenue projections were systematically optimistic
\item \textbf{Risk underestimation}: The probability of adverse economic shocks was not adequately factored into decisions
\item \textbf{Political agency problems}: Legislators could claim credit for projects while deferring the fiscal costs to future administrations
\end{enumerate}

The lesson that states drew in the 1840s was that ``taxes must be raised when spending is contemplated.'' If current taxes are not raised, taxpayers and politicians may not adequately factor in the risks of higher taxes in the future \citep{grinath1997debt}.

\subsection{Maryland and Pennsylvania: The Critical Difference}

\citet{wallis2004sovereign} highlights a critical difference between states that defaulted and states with comparable debt levels that did not. Maryland and Pennsylvania both had the fiscal capacity to service their debts through property taxation. However, both states were ``dilatory in levying adequate taxes.''

In contrast, when faced with fiscal pressures, Maryland and Pennsylvania eventually enacted property taxes and were able to resume debt payments with back interest as soon as these taxes were implemented. This demonstrates that the problem was not inability to pay, but unwillingness to tax until crisis forced action.

The northwestern states (Indiana, Illinois, Michigan), by contrast, had already implemented high property taxes and genuinely ran out of resources to tax. These states were forced to default and then renegotiate or partially repudiate their debts.

\section{Long-Term Impact}

The constitutional restrictions adopted in the 1840s and 1850s proved remarkably durable and effective. No state has experienced a full default since Arkansas in 1933, despite numerous economic downturns and fiscal pressures over the intervening 90+ years.

The reforms established several enduring principles in American state public finance:

\begin{itemize}
\item \textbf{Balanced budget requirements}: Most state constitutions now require balanced operating budgets
\item \textbf{Debt limitations}: States maintain various forms of constitutional or statutory debt limits
\item \textbf{Voter approval for borrowing}: Many states require voter approval for general obligation bonds
\item \textbf{Separation of state and private enterprise}: States generally cannot invest in or lend to private companies
\end{itemize}

These restrictions have been criticized at times for limiting states' flexibility in responding to economic conditions or funding needed infrastructure. However, they represent a deliberate choice, made in the wake of fiscal crisis, to constrain the ability of political actors to defer fiscal burdens to the future.

\section{Conclusion}

The state defaults of the 1840s and the constitutional responses they provoked represent a critical episode in the development of American fiscal federalism. States learned, through bitter experience, that infrastructure investment funded through debt without corresponding taxation created systematic risks that voters and legislators were ill-equipped to assess.

The constitutional reforms adopted by New York, Pennsylvania, and other states addressed this problem by eliminating taxless finance. By requiring that taxes be approved at the time debt is incurred, and by prohibiting state involvement in private enterprise, these reforms fundamentally altered the political economy of state borrowing.

These institutions have proven remarkably stable and effective. The fiscal constraints imposed in the wake of the 1840s crisis continue to shape state fiscal policy nearly two centuries later, standing as testament to the power of constitutional rules to constrain political behavior and prevent the recurrence of past mistakes.

\bibliographystyle{aer}
\bibliography{state_defaults}

\end{document}
